\begin{abstract}
Multiple-choice scores are commonly interpreted relative to nominal chance, even when an input-blind predictor can do better. We separate two reporting questions. For interface comparison, an arm-independent reference uses the construct's best constant input-blind predictor. For item-level information in one arm, an arm-conditional permutation null preserves the arm's prediction marginal within option-count strata. The first reference is a one-directional reporting check, not a certificate of construct validity; the second assigns every pure constant emitter exactly zero regardless of its collapse letter.

The best-constant floor must itself be calibrated under a null that the construct can realise. Restricting that null to each item's legal letters leaves only MMLU and BoolQ with label floors that a balanced null could not plausibly produce; for the other letter constructs, the analysis motivates a different reference without establishing a floor--chance difference in magnitude. On MMLU-Pro, only 1/12 non-OLMo cells clears the best-constant floor under the matched paired-bootstrap standard. Across five smaller benchmarks, 9/85 designated damaged cells clear the floor on the interval standard, all of them healed OLMo-2 arms; achieved precision explains why most per-benchmark tests remain inconclusive (Table~\ref{tab:headline-reference}).

Applying the arm-conditional null to 27 cells withdraws one above-floor capability label, removes both below-floor competence labels, and exposes a weak intact-family anchor. Full-fp32 evaluation eliminates every bf16 exact tie without recovering accuracy, so exact numerical ties do not explain the observed failure. These results support a practical rule: state the construct-specific input-blind reference, its uncertainty, and its intended inferential role before interpreting arm differences.
\end{abstract}
